\section{A causal transfer theorem with retained state}\label{sec:theorem}
\subsection{Raw interface and one common program}
Let $\Theta$ be a nonempty set. A parameter $\theta\in\Theta$ is fixed before the experiment and is not readable by the observer. Physical states, reports and performed audit marks are standard Borel. There is a finite ordinary action set $A$. A positive normalized marked kernel specifies the next physical state, the visible report, the visible end-of-cycle flag and the actual audit mark. Costs belong to $[0,1]$. Audit marks are not feedback unless explicitly included in the visible report. Decisions belong to a compact metric space $\mathcal D$, and the bounded marked loss is continuous in the decision, uniformly in the other variables when a quantitative modulus is used.

A visible cycle boundary triggers a prescribed, charged physical preparation. Preparation is call one of the next cycle, and its first ordinary data report is at call two. The physical preparation has a fixed law conditional on $\theta$, independent of previous innovations. The observer label is \emph{not overwritten}. A known preparation-report row, possibly the identity, is part of the program. No extra intercycle seed, tape, counter or exact posterior register is available. All internal randomization is freshly drawn conditional on the stored label. At a scored call the new label is committed before its fixed readout is used; a cleared report workspace is not a second, uncounted readout channel. For simplicity each cycle contains at least one ordinary call; allowing a preparation-only cycle changes no argument.

A stationary $M$-state program $\gam$ consists of action probabilities $\alpha_i(a)$, Borel report rows $k_{ia}(y,j)$, the preparation row, and readouts $v_i\in\mathcal D$, where $i,j\in[M]$. These same rows are used at every occurrence, for every $\theta$. Visible report types may distinguish preparation, ordinary and end reports; they are inputs, not uncounted persistent state. Every ordinary action is legal at every ordinary call. Other finite phase interfaces can be encoded in the $M$ labels and are counted there. We write $\Gamma_M$ for this common program space and start from the fixed label $1$.

The physical fresh-start property makes the complete cycle law conditional on an entering label independent of earlier history. For fixed $(\theta,\gam)$ write
\begin{equation}\label{eq:cycle-data}
 P_{ij}=\PP_i(I_{\rm next}=j),\qquad
 d_i=\E_i\tau,\qquad r_i=\E_i R,
 \qquad 1\le d_i\le\bar\mu,\quad 0\le r_i\le d_i.
\end{equation}
Here $\tau$ includes preparation, every failed acquisition, and the ending call; $R$ is their total bounded loss. The triples $(\tau,R,I_{\rm next})$ need not be iid across cycles. They form a Markov-renewal kernel on the actual retained labels. Dependence between duration, reward and next label is allowed.

We impose two conditions before optimizing a program.

\smallskip\noindent\textbf{Finite-response domination.}
For each action $a$ fix a probability reference $\sigma_a$ on its visible report space, with visible types kept distinct. For every finite intervention word and every entering label, the report law and its bounded performed-payoff submeasures have $L^1$ densities relative to the corresponding finite product of these references. The densities are raw experiment data, not chosen for each controller. A preparation coordinate has its own common reference. Zeros and changing support inside a reference are allowed. Domination of single-report marginals alone is not sufficient.

\smallskip\noindent\textbf{Uniform physical return.}
There is a displayed decreasing function $a:\N_0\to[0,\infty)$ such that
\begin{equation}\label{eq:tail}
 \sup_{\theta,\gam\in\Gamma_M,i}\E_i^{\theta,\gam}(\tau-T)_+
 \le a(T),\qquad a(T)\longrightarrow0,
 \qquad\bar\mu=a(0)<\infty.
\end{equation}
The supremum can instead range over a specified closed raw-legal architecture; that architecture must then replace $\Gamma_M$ everywhere. In particular the condition is conditional on \emph{every} entering label. Section~\ref{sec:raw} derives it from an action-uniform killed-kernel drift, not an independent clock assumption.

Give the ordinary and preparation report rows their weak-star topology in the respective reference $L^\infty$ spaces, action rows their finite simplex topology, and readouts their compact topology. A prescribed preparation row, such as the identity, is a fixed element of that space. Equalities are taken modulo the fixed reference null sets. This is a compact program space.

\subsection{The boundary regularity quantity}
For any finite stochastic matrix $P$ define
\begin{equation}\label{eq:group}
 \Pi_P=\lim_{n\to\infty}{1\over n}\sum_{k=0}^{n-1}P^k,
 \qquad G_P=(I-P+\Pi_P)^{-1}-\Pi_P,
 \qquad h(P)=\norm{G_P}_\infty.
\end{equation}
The matrix norm is the maximum absolute row sum. The limit and inverse exist for every finite stochastic matrix, including reducible and periodic ones. This is the group inverse of $I-P$ \cite{Meyer}. Fix $H\ge0$ and set
\begin{equation}\label{eq:regularclass}
 \Gamma_{M,H}^{\Theta}
 =\{\gam\in\Gamma_M:h(P_\theta^\gam)\le H\text{ for every }\theta\in\Theta\}.
\end{equation}
This is a restriction on the induced recurrence, not an additional memory register. Its nonemptiness is assumed whenever its value is discussed. All Borel controllers satisfying the displayed restriction are included.

For each cycle matrix put
\begin{equation}\label{eq:timechange}
 D=\diag(d_i),\qquad S=I-D^{-1}(I-P),\qquad c=D^{-1}r,
 \qquad g_\theta(\gam)=e_1\Pi_Sc.
\end{equation}
The matrix $S$ is stochastic because $d_i\ge1$. If $C$ ranges over the closed classes of $P$, $\pi_C$ is its stationary law on $C$, and $w_C$ is the probability of eventually entering $C$ from label $1$, then
\begin{equation}\label{eq:classratio}
 g_\theta(\gam)=\sum_C w_C\,{\pi_C r\over\pi_C d}.
\end{equation}
The individual ratios are taken \emph{before} mixing classes. In general $(e_1\Pi_Pr)/(e_1\Pi_Pd)$ is a different, incorrect average.

Let $C_N$ be the loss over the first $N$ physical calls, starting at preparation, and set
\[
 J_N^\theta(\gam)=\E C_N/N,
 \qquad J_\beta^\theta(\gam)=(1-\beta)\E\sum_{t\ge1}\beta^{t-1}\ell_t.
\]
The evaluation index $N$ and the discount factor are not stored in the running machine. Write $V_N^\Theta(M,H)$, $V_\beta^\Theta(M,H)$ and $V_\av^\Theta(M,H)$ for the infimum of the respective worst-model losses over the \emph{same} class \eqref{eq:regularclass}.

\begin{theorem}[Causal transfer with retained boundary state]\label{thm:main}
Under finite-response domination and \eqref{eq:tail}, the following hold for every $M\ge1$ and every nonempty $\Gamma_{M,H}^{\Theta}$.

\textup{(i)} The class \eqref{eq:regularclass} is compact. The average loss exists almost surely for each fixed model and program, has expectation \eqref{eq:classratio}, and
\begin{equation}\label{eq:attain}
 V_\av^\Theta(M,H)=\min_{\gam\in\Gamma_{M,H}^\Theta}\sup_{\theta\in\Theta}g_\theta(\gam)
 =\sup_{\substack{F\subset\Theta\\0<|F|<\infty}} V_\av^F(M,H).
\end{equation}
The class on the right imposes the recurrence bound only for the models in $F$. A risk level strictly below the left side is excluded by a finite family of joint risk-and-recurrence constraints.

\textup{(ii)} Put $B=(2M+1)H$ and
\begin{equation}\label{eq:bN}
 b_N=\min\left\{1,{3\bar\mu B+\sum_{s=0}^Na(s)\over N}\right\}.
\end{equation}
Uniformly over $\theta$ and the declared programs,
\begin{equation}\label{eq:bridges}
 |J_N^\theta-g_\theta|\le b_N,\qquad
 |J_\beta^\theta-g_\theta|
 \le(1-\beta)^2\sum_{N\ge1}N\beta^{N-1}b_N.
\end{equation}
The same bounds hold between their robust optimal values. Both errors tend to zero. Every cluster point of asymptotically optimal finite-acquisition or discounted programs is average-optimal in \eqref{eq:regularclass}.

\textup{(iii)} If $\sup_{\theta,\gam,i}\E_i\tau^{1+p}\le K$ for $0<p\le1$, the respective errors in \eqref{eq:bridges} are bounded by
\begin{equation}\label{eq:moment-rate}
 {3\bar\mu B\over N}+2^{1-p}MKN^{-p},\qquad
 3\bar\mu B(1-\beta)+2^{1-p}MK(1-\beta)^p,
\end{equation}
capped at one. Boundary recurrence and physical return moments are distinct quantities.

\textup{(iv)} For a finite model family, let $V_\av^\Theta(M)$ be the infimum over \emph{all} programs in $\Gamma_M$, with no recurrence bound. Then
\begin{equation}\label{eq:exhaustion}
 V_\av^\Theta(M)=\inf_{k\in\N}V_\av^\Theta(M,k).
\end{equation}
Unrestricted attainment is equivalent to the infimum in \eqref{eq:exhaustion} being reached at some nonempty integer level. No uniform finite-$N$ or discount interchange is asserted after taking this unbounded exhaustion.

\textup{(v)} For a specified squared-prediction task, the actual class-weighted acquired law and the counted controller--predictor construction of Theorem~\ref{thm:geometry} give the joint lower and upper curves \eqref{eq:optimizedlower}--\eqref{eq:optimizedupper}. Under the additional uniform mass, global-cover, candidate-stability and attaining-exploration hypotheses of Corollary~\ref{cor:matching}, they imply
\begin{equation}\label{eq:centralgeometry}
 cM^{-2/d}-b_N\ \le\ V_N(M,H)-B_*\ \le\ CM^{-2/d}+b_N,
\end{equation}
with constants independent of $M,N$ in that declared class. The same-memory numerical root bound is \eqref{eq:numericalgeom}; a complete-cycle simulation adds the explicit modulus \eqref{eq:defectmod} at the \emph{joint} simulator recurrence bound. Arbitrary allowed error tolerances do not give lower-risk terms; Theorem~\ref{thm:joint} provides an attained joint memory/calibration/defect law. Thus the geometric part is a conditional matching consequence of the same causal transfer, not a dimension law for every raw kernel.
\end{theorem}

The theorem includes genuine retained information, but its uniform statement is not over all possible boundary degenerations. The next theorem describes that boundary exactly for cost-universal Cesaro regularity. For an infinite model family an individual program can have $\sup_\theta h(P_\theta)=\infty$; \eqref{eq:exhaustion} is therefore stated only for finite families.
